\section{Exact Teaching Prompts and Example Responses}
\label{app:materials-examples}

\subsection{Teaching Prompts}

Every main system message starts with one of the following strings. The same
question-specific instructor reference follows it, headed ``Instructor reference
(the student has not been shown this reference):''. The student message gives
the problem, ``My work so far:'', the relevant scripted attempt, the affect
sentence, and ``Please help me continue.'' The two affect sentences are
``I feel calm while working on this.'' and ``I feel frustrated while working
on this.'' Full materials and input hashes accompany the generation manifests.

\paragraph{Standard prompt (Canonical in tables).}
``You are a tutor. Respond to the student's last message.''

\paragraph{Neutral A.}
``You are helping a student as a tutor. Reply to the student's most recent message.''

\paragraph{Neutral B.}
``Act as a tutor assisting the student. Write your response to their latest message.''

\paragraph{ASK.}
The standard prompt followed by:
``For this response, prioritize asking the student to work out the next useful
step. Do not state the final answer. You may offer a brief hint or explanation.''

\paragraph{EXPLAIN.}
The standard prompt followed by:
``For this response, prioritize explaining the solution and state the final
answer. You may invite follow-up reasoning or questions.''

\subsection{One Input, Five Responses}

For illustration, we use the first template in the finalized test set,
with the standard tutoring prompt, erroneous student work, the frustrated
cue, and the first request.
This example set was selected after the analyses and is not used to estimate
frequency or validate the coding rubric. It retains all five models,
including a disputed elicitation code.

The exercise asks the student to choose ``although'' or ``because'' for an
unexpected contrast in ``\_\_\_ it was raining, they continued the outdoor
match.'' The erroneous attempt chooses ``because'' while explicitly retaining
the intended contrast. All five responses disclose ``although'', including
GLM, despite its much lower pooled answer-giving rate under the standard prompt. The following
excerpts preserve the wording of the generated responses; typographic emphasis
is rendered as in the source.

\paragraph{Generic praise versus acknowledgement.}
M3 begins, ``You're absolutely right to question your choice!'' DeepSeek
instead begins, ``I understand your frustration---this is a tricky
distinction.'' Doubao begins, ``It's totally okay to feel frustrated here''.
The complete M3 response is coded as not acknowledging affect/difficulty,
while the other two are coded positive by all three annotators. The contrast
illustrates the intended separation between general encouragement and explicit
acknowledgement; unanimous coding is still not correct interpretation of the behavior.

\paragraph{A disclosed answer and a further invitation.}
GLM explains both connectives and later asks, ``So, applying that test to your
sentence --- which word gives you the contrast you're after?'' Its revelation
and elicitation codes are both unanimously positive. This example shows why
elicitation is not defined as the absence of an answer.

\paragraph{A coding boundary.}
M2.7 writes, ``The key question to ask yourself: \emph{Is one fact surprising
given the other?} If yes, reach for \textbf{although}.'' Its elicitation code
is majority-positive but not unanimous. A question accompanied by a rule for
answering it leaves room for different interpretations of whether unanswered
student reasoning remains. We retain this disagreement; the illustration does
not override the fixed majority rule or supply a new reference label.

\begin{table}[htbp]\centering\small
\begin{tabular}{lrrr}\toprule
Model & Reveal & Elicit & Acknowledge \\
\midrule
M2.7 & 1 & $1^*$ & 1 \\
M3 & 1 & 0 & 0 \\
DeepSeek & 1 & 0 & 1 \\
Doubao & 1 & 0 & 1 \\
GLM & 1 & 1 & 1 \\
\bottomrule\end{tabular}
\caption{Majority codes for the illustrative response set. All cells are
unanimous except the starred M2.7 elicitation code. This single input is not
representative evidence about model frequency, and no label is revised for the illustration.}
\label{tab:illustrative-codes}\end{table}
